%*******************************************************
% 4.3.7 đến 4.3.9
%*******************************************************
\subsection{Kiểm định giả thuyết}\label{sec:kiemdinh}

\begin{tomtat}
\HH{3} được kiểm định bằng khác biệt kép: so thay đổi thanh khoản của cổ phiếu chuyển sàn với thay đổi của cổ phiếu đối chứng. Vì mọi cổ phiếu chuyển sàn cùng một ngày, suy luận không dựa vào sai số chuẩn thông thường mà dựa vào bootstrap tự hồi quy trên chuỗi chênh lệch theo tháng. Phiên bản nghiên cứu sự kiện kiểm tra xu hướng song song. Các tương tác cho biết tác động khác nhau thế nào theo mức cổ phiếu chịu ba thay đổi quy tắc (biên độ, bước giá, việc vào chỉ số HOSE) và theo trụ sở doanh nghiệp. \HH{4} được báo cáo bằng ước lượng; \HH{1} được kiểm định trên dữ liệu bảng quốc tế.
\end{tomtat}

\subsubsection*{Phương trình khác biệt kép}

Đề tài ước lượng tác động của việc chuyển sàn bằng khác biệt kép trên dữ liệu bảng cổ phiếu-tháng (\HH{3}). Với cổ phiếu $i$ và tháng $t$, phương trình khác biệt kép là:
\begin{equation}\label{eq:did}
A_{it} = \alpha_i + \tau_t + \delta D_i\Post_t + \varepsilon_{it}
\end{equation}
trong đó $A_{it}$ là logarit của độ phi thanh khoản Amihud của cổ phiếu $i$ trong tháng $t$, tính trên các ngày có lợi suất tuyệt đối dưới ngưỡng chung và không liền sau một ngày có lợi suất tuyệt đối từ ngưỡng này trở lên (mục~\ref{sec:giathuyet}); $D_i$ bằng 1 nếu cổ phiếu chuyển sàn, tức mọi cổ phiếu niêm yết tại HNX được chuyển sang HOSE theo lịch của VNX, không chỉ các mã của rổ CS; $\Post_t$ bằng 1 từ tháng 1/2027; $\alpha_i$ và $\tau_t$ là hiệu ứng cố định theo cổ phiếu và theo tháng; $\delta$ là tác động của việc chuyển sàn; $\varepsilon_{it}$ là sai số.

\paragraph{Đọc công thức \eqref{eq:did} bằng lời.} Thanh khoản của một cổ phiếu trong một tháng gồm bốn phần: một mức riêng của cổ phiếu ấy, không đổi theo thời gian ($\alpha_i$); một cú sốc chung của tháng ấy, tác động lên mọi cổ phiếu ($\tau_t$); phần thay đổi thêm chỉ có ở cổ phiếu chuyển sàn sau ngày chuyển ($\delta$, vì tích $D_i\Post_t$ chỉ bằng 1 khi cả hai điều kiện cùng đúng); và phần còn lại không giải thích được ($\varepsilon_{it}$). Tác động của việc chuyển sàn là $\delta$.

Vì $A_{it}$ là logarit, khi $\delta$ nhỏ thì $100\delta$ xấp xỉ phần trăm thay đổi; chính xác hơn, độ phi thanh khoản thay đổi $(e^{\delta} - 1)$ nhân 100 phần trăm. \HH{3} kỳ vọng $\delta$ âm, tức cổ phiếu chuyển sàn dễ mua bán hơn.

\begin{vidu}[: khác biệt kép trên bốn con số]
Giả sử bình quân logarit Amihud như sau:
\begin{center}
\begin{tabular}{@{}lccc@{}}
\toprule
 & Trước sự kiện & Sau sự kiện & Thay đổi\\
\midrule
Cổ phiếu chuyển sàn & 1,80 & 1,40 & $-0{,}40$\\
Cổ phiếu đối chứng & 1,20 & 1,05 & $-0{,}15$\\
\midrule
Khác biệt kép & & & $-0{,}25$\\
\bottomrule
\end{tabular}
\end{center}
Cổ phiếu chuyển sàn dễ mua bán hơn ($-0{,}40$), nhưng cổ phiếu đối chứng cũng dễ mua bán hơn ($-0{,}15$), có lẽ vì cả thị trường tốt lên. Phần chỉ do chuyển sàn là $-0{,}40 - (-0{,}15) = -0{,}25$. Quy ra phần trăm: $e^{-0{,}25} - 1 \approx -0{,}22$, tức độ phi thanh khoản giảm khoảng 22\%. Nếu chỉ so trước và sau ở nhóm chuyển sàn, ta sẽ nói nhầm là giảm khoảng 33\% ($e^{-0{,}40} - 1$).
\end{vidu}

\begin{vidao}
\emph{Vì sao dùng khác biệt kép.} Có hai cách đơn giản hơn, và cả hai đều sai theo một cách dễ thấy. So cổ phiếu chuyển sàn trước và sau sự kiện thì lẫn tác động của sàn với mọi thay đổi chung của thị trường trong cùng thời gian (lãi suất, dòng vốn ngoại, nâng hạng). So cổ phiếu chuyển sàn với cổ phiếu HOSE sau sự kiện thì lẫn tác động của sàn với mọi khác biệt sẵn có giữa hai nhóm doanh nghiệp. Khác biệt kép trừ đi cả hai: hiệu ứng cố định theo cổ phiếu loại khác biệt sẵn có, hiệu ứng cố định theo tháng loại thay đổi chung. Cái giá phải trả là giả định xu hướng song song (mục~\ref{sec:did}).
\end{vidao}

Phương trình \eqref{eq:did} được ước lượng trên giai đoạn tham chiếu (đến tháng 6/2025) và giai đoạn sau sự kiện (từ tháng 1/2027); các tháng từ 7/2025 đến 12/2026, gồm tháng chuyển sàn, chỉ được đưa vào phương trình \eqref{eq:es}. Vì giai đoạn tham chiếu dài 16 năm, cửa sổ trước của ước lượng chính bắt đầu sau điểm gãy cấu trúc cuối cùng (Bai \& Perron, 1998) của chuỗi chênh lệch logarit Amihud theo tháng giữa cổ phiếu chuyển sàn và cổ phiếu đối chứng (chuỗi $d_t$, định nghĩa đầy đủ ở phần suy luận bên dưới), và ước lượng với cửa sổ trước 24 tháng (7/2023--6/2025) được báo cáo kèm.

Lý do chọn cửa sổ trước theo điểm gãy cấu trúc là thị trường năm 2010 khác xa thị trường năm 2024. Nếu chênh lệch giữa hai nhóm đã đổi mức ở một thời điểm trong quá khứ, chẳng hạn khi quy chế giao dịch thay đổi, thì dùng toàn bộ 16 năm làm \enquote{trước sự kiện} sẽ trộn hai chế độ khác nhau. Bắt đầu từ sau điểm gãy cuối cùng giúp cửa sổ trước chủ yếu gồm giai đoạn có cùng chế độ với thời điểm ngay trước sự kiện; vì ngày gãy được ước lượng có sai số, ước lượng với cửa sổ trước 24 tháng được báo cáo kèm.

\subsubsection*{Nhóm đối chứng}

Nhóm đối chứng chính gồm cổ phiếu HOSE ghép láng giềng gần nhất theo vốn hóa, ngành và độ phi thanh khoản tại ngày 30/6/2025, có hoàn lại, trong vùng hỗ trợ chung. Vì ghép theo độ phi thanh khoản có thể tạo hồi quy về trung bình, kết quả khi chỉ ghép theo vốn hóa và ngành được báo cáo thêm (Daw \& Hatfield, 2018). Nhóm đối chứng phụ gồm cổ phiếu UPCoM trong thời gian thị trường này còn do HNX vận hành.

\enquote{Vùng hỗ trợ chung} nghĩa là chỉ giữ những cổ phiếu chuyển sàn tìm được cổ phiếu HOSE đủ giống; mã chuyển sàn quá nhỏ hay quá kém thanh khoản, không có cặp tương xứng ở HOSE, bị loại khỏi phép so sánh chính.

\begin{delam}
\emph{Dễ nghĩ là} ghép càng nhiều tiêu chí thì càng tốt, kể cả ghép theo chính thước đo kết quả. Nhưng nếu một cổ phiếu HOSE được chọn vì tháng 6/2025 nó có độ phi thanh khoản cao bất thường, thì các tháng sau nó có xu hướng trở về mức bình thường của nó, tức thanh khoản \enquote{tốt lên}. Nhóm đối chứng tốt lên do quy luật thống kê này làm $\delta$ bị kéo về 0. Daw và Hatfield (2018) cảnh báo đúng hiện tượng này; đề tài vì vậy báo cáo thêm kết quả khi không ghép theo độ phi thanh khoản.
\end{delam}

\subsubsection*{Suy luận khi chỉ có một sự kiện}

Vì mọi cổ phiếu chuyển sàn cùng một ngày, suy luận chính dựa trên chuỗi $d_t$, là chênh lệch trung bình logarit Amihud giữa cổ phiếu chuyển sàn và cổ phiếu đối chứng ghép cặp trong tháng $t$. Thống kê kiểm định là trung bình của $d_t$ sau sự kiện trừ trung bình trong cửa sổ trước. Giá trị p chính được tính bằng bootstrap tự hồi quy trên $d_t$, với mô hình tự hồi quy ước lượng trên giai đoạn tham chiếu, hệ số được hiệu chỉnh độ chệch, và kích thước kiểm định (xác suất bác bỏ giả thuyết gốc khi thực ra không có tác động) được kiểm tra bằng mô phỏng trước ngày đăng ký.

\begin{vidao}
\emph{Vì sao không dùng sai số chuẩn phân cụm theo cổ phiếu.} Sai số chuẩn phân cụm theo cổ phiếu giả định rằng các cổ phiếu là những quan sát độc lập. Nhưng ở đây, mọi cổ phiếu chuyển sàn cùng chịu một sự kiện vào cùng một ngày. Cú sốc tác động như nhau lên mọi cổ phiếu đã được hiệu ứng cố định theo tháng $\tau_t$ hấp thụ; vấn đề là những cú sốc theo tháng tác động lên cả nhóm chuyển sàn khác với nhóm đối chứng (một tin chỉ liên quan đến các mã HNX, một đợt bán tháo dồn vào mã nhỏ), làm sai số của các mã chuyển sàn tương quan với nhau trong cùng tháng. Có hàng trăm cổ phiếu nhưng về bản chất chỉ có \emph{một} lần chuyển sàn. Nếu coi 299 mã là 299 thí nghiệm, sai số chuẩn sẽ nhỏ đến mức gần như mọi khác biệt đều \enquote{có ý nghĩa}. Conley và Taber (2011) chỉ ra vấn đề này cho khác biệt kép với ít lần thay đổi chính sách.

Cách làm của đề tài là gộp toàn bộ thông tin thành một chuỗi theo tháng, $d_t$, rồi hỏi: nếu không có sự kiện, chuỗi này tự dao động lớn đến đâu? Mô hình tự hồi quy mô tả cách $d_t$ dao động trong giai đoạn tham chiếu, kể cả việc tháng này giống tháng trước. Bootstrap tạo ra hàng nghìn chuỗi giả theo đúng quy luật ấy, tính thống kê kiểm định trên mỗi chuỗi, và giá trị p là tỷ lệ chuỗi giả cho thống kê âm hơn thống kê thật.
\end{vidao}

\begin{vidu}[: bootstrap tự hồi quy trong năm bước]
\begin{enumerate}
\item Tính $d_t$ cho mọi tháng của cửa sổ trước và cửa sổ sau.
\item Trên giai đoạn tham chiếu, ước lượng mô hình tự hồi quy, ở đây minh họa bằng bậc một: $d_t = c + \phi d_{t-1} + u_t$; giả sử được $\phi = 0{,}6$ sau khi hiệu chỉnh độ chệch, và giữ lại các phần dư $u_t$ đã trừ trung bình.
\item Tạo một chuỗi giả theo đúng mô hình đã ước lượng, $d^*_t = c + \phi d^*_{t-1} + u^*_t$, với $u^*_t$ rút ngẫu nhiên có hoàn lại từ các phần dư, bắt đầu từ một giá trị ban đầu và bỏ một đoạn khởi động. Chuỗi giả trải liên tục qua cửa sổ trước, 18 tháng chuyển tiếp và cửa sổ sau; các tháng chuyển tiếp không vào thống kê nhưng được giữ để tự tương quan giữa hai cửa sổ được tính đúng. Chuỗi giả này \enquote{không có sự kiện}.
\item Trên chuỗi giả, tính trung bình sau trừ trung bình trước.
\item Lặp lại bước 3--4 chẳng hạn 9.999 lần. Nếu thống kê thật là $-0{,}25$ và chỉ 120 trong 9.999 chuỗi giả cho giá trị từ $-0{,}25$ trở xuống, thì giá trị p một phía xấp xỉ 0,012.
\end{enumerate}
Số lần lặp và các con số trong ví dụ là minh họa; thủ tục cụ thể được ghi trong tệp đăng ký.
\end{vidu}

Suy luận trên các ngày sự kiện giả trong giai đoạn tham chiếu, theo tinh thần của Conley và Taber (2011), chỉ là kiểm tra bổ sung: các cửa sổ giả chồng lấn nhau, chỉ tạo khoảng sáu cửa sổ độc lập, và mô phỏng sơ bộ của thí sinh cho thấy kiểm định này bác bỏ giả thuyết gốc với tần suất cao hơn đáng kể so với mức danh nghĩa. \HH{3} được ủng hộ khi $\delta$ âm và giá trị p bootstrap tự hồi quy nhỏ hơn ngưỡng Holm tương ứng.

Ý tưởng của ngày sự kiện giả là giả vờ sự kiện xảy ra ở một tháng nào đó trong quá khứ, tính lại thống kê, và xem thống kê thật có cực đoan hơn phần lớn các thống kê giả không. Vấn đề là với cửa sổ sau dài 24 tháng, các ngày giả cách nhau vài tháng cho các cửa sổ gần như trùng nhau; số phép so sánh thật sự độc lập ít hơn nhiều so với số ngày giả, nên kiểm định này dễ bác bỏ sai.

\subsubsection*{Phiên bản nghiên cứu sự kiện}

Phiên bản nghiên cứu sự kiện kiểm tra xu hướng song song và phản ứng trong giai đoạn chuyển tiếp:
\begin{equation}\label{eq:es}
A_{it} = \alpha_i + \tau_t + \sum_k \delta_k D_i T_{kt} + \varepsilon_{it}
\end{equation}
trong đó $T_{kt}$ bằng 1 nếu tháng $t$ là tháng $k$, với $k$ chạy từ tháng 1/2024 đến tháng 12/2028 và bỏ tháng 6/2025 làm tháng gốc; các tháng trước 1/2024 được gộp thành một nhóm. Các hệ số $\delta_k$ trong giai đoạn chuyển tiếp cho biết phản ứng trước sự kiện. Đề tài kiểm định đồng thời giả thuyết các hệ số $\delta_k$ từ tháng 1/2024 đến tháng 5/2025 bằng 0. Độ nhạy của kết luận trước các vi phạm xu hướng song song được đánh giá theo Rambachan và Roth (2023). Vì giữa tháng gốc và giai đoạn sau có 18 tháng chuyển tiếp, khoảng tin cậy theo cách này có thể rất rộng.

\paragraph{Đọc công thức \eqref{eq:es} bằng lời.} Thay vì một hệ số $\delta$ duy nhất cho toàn bộ giai đoạn sau, mỗi tháng có một hệ số riêng $\delta_k$, đo chênh lệch giữa cổ phiếu chuyển sàn và đối chứng trong tháng $k$ so với tháng gốc 6/2025. Vẽ các $\delta_k$ theo thời gian sẽ cho một đồ thị: trước 7/2025 các điểm nên nằm quanh 0 (xu hướng song song), trong giai đoạn chuyển tiếp có thể bắt đầu lệch (phản ứng sớm), sau 1/2027 cho thấy tác động diễn biến ra sao (tăng dần, ổn định hay mờ đi như ở Tokyo).

\begin{vidu}
Giả sử các $\delta_k$ từ 1/2024 đến 5/2025 dao động trong khoảng $-0{,}03$ đến $0{,}04$ và kiểm định đồng thời không bác bỏ giả thuyết chúng bằng 0: đây là dấu hiệu ủng hộ xu hướng song song. Nếu ngược lại, các $\delta_k$ giảm đều từ $0$ xuống $-0{,}15$ trong cùng giai đoạn, thì nhóm chuyển sàn đã có xu hướng tốt lên riêng từ trước, và một phần $\delta$ ước lượng được sẽ không phải do chuyển sàn. Phương pháp của Rambachan và Roth (2023) cho biết kết luận còn đứng vững không nếu xu hướng riêng ấy kéo dài sang giai đoạn sau với độ lớn tương tự.
\end{vidu}

\subsubsection*{Tác động khác nhau theo mức chịu cú sốc}

Việc chuyển sàn đổi nhiều quy tắc cùng lúc, trong đó ba quy tắc tác động mạnh yếu khác nhau tùy cổ phiếu: biên độ hẹp lại chỉ ràng buộc mã thường có lợi suất ngày lớn; bước giá nhỏ đi chủ yếu ở mã giá thấp; việc vào VN-Index và các chỉ số khác của HOSE tạo áp lực cầu lớn hơn ở mã có vốn hóa lớn so với giá trị giao dịch (mục~\ref{sec:did}).

Ba mức chịu cú sốc tương ứng được đo bằng dữ liệu đến ngày 30/6/2025: tỷ lệ ngày có lợi suất tuyệt đối từ 7\% đến dưới 10\%; nghịch đảo của giá đóng cửa bình quân; và tỷ số giữa vốn hóa toàn phần và giá trị giao dịch bình quân ngày, vì VN-Index tính trọng số theo vốn hóa toàn phần, không điều chỉnh tỷ lệ tự do chuyển nhượng (Sở Giao dịch Chứng khoán Thành phố Hồ Chí Minh, 2026).

\begin{center}
\captionof{table}{Ba thay đổi quy tắc khi chuyển sàn và thước đo mức chịu cú sốc}\label{tab:cusoc}
{\bangnho
\begin{tabularx}{\textwidth}{@{}>{\raggedright\arraybackslash}p{0.2\textwidth} L L@{}}
\toprule
Cú sốc & Cổ phiếu chịu mạnh nhất & Thước đo mức chịu cú sốc\\
\midrule
Biên độ hẹp lại & Mã thường có ngày tăng giảm 7--10\% & Tỷ lệ ngày có lợi suất tuyệt đối từ 7\% đến dưới 10\%\\
Bước giá nhỏ đi & Mã giá thấp & Nghịch đảo giá đóng cửa bình quân\\
Vào chỉ số HOSE & Mã vốn hóa lớn so với giá trị giao dịch & Vốn hóa toàn phần chia giá trị giao dịch bình quân ngày\\
\bottomrule
\end{tabularx}}
\end{center}

Bộ quy tắc của FTSE Russell coi cả HOSE và HNX là sàn đủ điều kiện ở Việt Nam (FTSE Russell, 2026a), nên việc đổi sàn tự nó không làm một mã mất hay có thêm tư cách đủ điều kiện và không được coi là một kênh tác động qua chỉ số FTSE. Danh sách cổ phiếu đủ điều kiện, lập với dữ liệu đến ngày 30/6/2026, chỉ gồm mã niêm yết tại HOSE (FTSE Russell, 2026b) và được dùng để loại mã khỏi nhóm đối chứng; nếu mã chuyển sàn nào được đưa vào chỉ số ở các kỳ sau, $\delta$ được ước lượng thêm khi loại các mã này.

Phương trình \eqref{eq:did} được mở rộng bằng các số hạng $D_i\Post_tE_{si}$, trong đó $E_{si}$ là mức chịu cú sốc $s$ của cổ phiếu $i$, trừ trung bình và chia cho độ lệch chuẩn giữa các mã chuyển sàn, nên $\delta$ là tác động tại mức chịu cú sốc trung bình. Các số hạng $E_{si}\tau_t$ cho mỗi mức chịu cú sốc một hiệu ứng theo tháng riêng, ước lượng trên mã đối chứng đo được cùng thước đo: mã UPCoM với biên độ, mã HOSE với bước giá và chỉ số. Nếu cổ phiếu UPCoM chuyển sang HOSE trước hoặc trong cửa sổ sau (mục~\ref{sec:tincay}), mức chịu cú sốc biên độ của mọi mã được đo bằng tỷ lệ ngày có lợi suất tuyệt đối từ ngưỡng chung trở lên và hiệu ứng theo tháng tương ứng được ước lượng trên mã HOSE đối chứng; thước đo này chỉ đại diện gần đúng cho mức chịu ràng buộc của biên độ.

\begin{vidu}[: đọc hệ số tương tác]
Giả sử $\delta = -0{,}20$ và hệ số của $D_i\Post_tE_{si}$ với cú sốc bước giá là $-0{,}08$. Vì $E_{si}$ đã được chuẩn hóa, một mã có mức chịu cú sốc bước giá ở đúng trung bình có tác động $-0{,}20$; một mã cao hơn trung bình một độ lệch chuẩn (giá thấp hơn) có tác động $-0{,}20 - 0{,}08 = -0{,}28$; một mã thấp hơn trung bình một độ lệch chuẩn có tác động $-0{,}12$. Kết quả như vậy gợi ý bước giá nhỏ đi là một kênh, nhưng như đoạn dưới giải thích, không tách được riêng tác động của bước giá.
\end{vidu}

Giả định nhận dạng là: nếu không chuyển sàn, chênh lệch thanh khoản giữa mã chịu cú sốc mạnh và mã chịu cú sốc yếu trong nhóm chuyển sàn sẽ biến động như ở nhóm đối chứng; giả định này được kiểm tra bằng các hệ số trước sự kiện. Vì các mức chịu cú sốc có thể tương quan với nhau và với quy mô, hệ số tương tác không đo riêng tác động của từng quy tắc; mô hình có thêm tương tác với quy mô và với logarit Amihud ban đầu.

Mỗi hệ số tương tác được suy luận như $\delta$, vì sai số chuẩn phân cụm theo cổ phiếu bác bỏ quá thường xuyên khi cú sốc chung theo tháng tác động mạnh hơn lên mã chịu cú sốc mạnh: với từng tháng, đề tài lấy hệ số góc của logarit Amihud theo mức chịu cú sốc ở mã chuyển sàn trừ hệ số góc ở mã đối chứng, rồi áp dụng bootstrap tự hồi quy cho chuỗi này. Các hệ số tương tác, kể cả theo trụ sở, là phân tích phụ đăng ký trước, không đưa vào thủ tục Holm.

\subsubsection*{Tác động khác nhau theo trụ sở doanh nghiệp}

Phân tích theo trụ sở kiểm tra xem việc chuyển sàn có tác động khác nhau đến doanh nghiệp có trụ sở tại Hà Nội và doanh nghiệp ở nơi khác hay không. Tài liệu cho thấy nhà quản lý quỹ ưa nắm giữ cổ phiếu của doanh nghiệp có trụ sở gần mình (Coval \& Moskowitz, 1999) và thanh khoản của cổ phiếu phụ thuộc vào vị trí trụ sở (Loughran \& Schultz, 2005), nhưng hai bài này nói về khoảng cách giữa nhà đầu tư và doanh nghiệp, không nói về vị trí của sàn. Kênh mà đề tài kiểm tra là sự chú ý và thông tin gắn với nơi đặt sàn: nếu kênh này có thật, cổ phiếu của doanh nghiệp có trụ sở tại Hà Nội được lợi ít hơn khi chuyển sang HOSE.

Chênh lệch này được ước lượng bằng tương tác giữa $D_i\Post_t$ và biến giả trụ sở tại Hà Nội, xác định theo địa chỉ đăng ký tại ngày 30/6/2025, kèm tích của biến giả này với hiệu ứng cố định theo tháng (ước lượng trên các mã HOSE đối chứng có trụ sở tại Hà Nội) và tương tác với quy mô và ngành.

Chênh lệch theo trụ sở là ước lượng khác biệt của tác động trong cùng một thiết kế, không phải tác động nhân quả của vị trí trụ sở, và không đo lượng thanh khoản mà thị trường vốn Hà Nội mất, vì mọi mã đều chuyển sang HOSE. Khoảng tin cậy hẹp quanh 0 sẽ không ủng hộ vai trò của vị trí sàn khi giao dịch đã điện tử hóa toàn quốc.

Kết quả này có ý nghĩa trực tiếp cho câu hỏi chính sách. Nếu khoảng tin cậy của chênh lệch theo trụ sở hẹp quanh 0, kết quả không ủng hộ vai trò của vị trí sàn đối với thanh khoản của các mã chuyển sàn; nếu khoảng tin cậy rộng, kết quả chưa đủ cơ sở kết luận. Trong cả hai trường hợp, chênh lệch này không đo lượng thanh khoản mà thị trường vốn Hà Nội mất. Nếu doanh nghiệp Hà Nội được lợi ít hơn rõ rệt, thì có một kênh địa phương đáng để thành phố quan tâm.

\subsubsection*{Báo cáo \HH{4}}

\HH{4} được báo cáo bằng chênh lệch của chỉ số khả năng đầu tư bình quân giữa giai đoạn sau sự kiện và giai đoạn tham chiếu, kèm khoảng tin cậy bootstrap khối dừng trên chuỗi lợi suất tuần của hai rổ, được lấy mẫu lại đồng thời, và số mô hình có cùng dấu. Chênh lệch tỷ số Sharpe, một thước đo phụ, được kiểm định theo Ledoit và Wolf (2008).

\enquote{Lấy mẫu lại đồng thời} nghĩa là mỗi lần bootstrap rút cùng các tuần cho cả rổ H và rổ CS. Hai rổ chịu chung cú sốc thị trường trong cùng tuần; rút riêng sẽ làm mất mối tương quan này và làm khoảng tin cậy của khoảng cách rộng hơn thực tế. \enquote{Số mô hình có cùng dấu} là một thước đo đơn giản về độ đồng thuận: nếu cả bảy mô hình đều cho khoảng cách thu hẹp, kết luận đáng tin hơn so với khi bốn mô hình thu hẹp còn ba mô hình nới rộng, dù bình quân có thể như nhau.

\subsubsection*{Kiểm định \HH{1}}

Đề tài kiểm định \HH{1} trên dữ liệu bảng của các trung tâm tài chính, dựa trên cách tiếp cận dữ liệu bảng của Naeem (2023), nhưng dùng điểm GFCI chuẩn hóa theo kỳ thay cho thứ hạng. Với trung tâm $i$ và kỳ công bố $t$, phương trình ước lượng là:
\begin{equation}\label{eq:gfci}
Z_{it} = \alpha_i + \tau_t + \beta\SMD_{it} + \lambda W_{it} + \varepsilon_{it}
\end{equation}
trong đó $Z_{it}$ là điểm GFCI của trung tâm $i$ ở kỳ $t$, chuẩn hóa thành điểm z trong từng kỳ vì điểm GFCI được ước lượng lại mỗi kỳ và số trung tâm thay đổi giữa các kỳ. $\SMD_{it}$ là chỉ báo phát triển thị trường cổ phiếu, gộp từ tỷ lệ vòng quay và vốn hóa trên GDP (Čihák và cộng sự, 2012), còn $W_{it}$ gồm thu nhập bình quân đầu người và chỉ số quản trị nhà nước (Kaufmann và cộng sự, 2011); cả hai được đo ở quốc gia chứa trung tâm $i$, trong năm liền trước kỳ $t$. $\lambda$ là véc-tơ hệ số của $W_{it}$; $\alpha_i$ và $\tau_t$ là hiệu ứng cố định theo trung tâm và theo kỳ; $\varepsilon_{it}$ là sai số. Nếu Cơ sở dữ liệu Phát triển Tài chính Toàn cầu không phủ các năm sau 2021, đề tài dùng số liệu của Liên đoàn các Sở Giao dịch Thế giới, sau khi kiểm tra độ phủ và điều khoản sử dụng.

\paragraph{Đọc công thức \eqref{eq:gfci} bằng lời.} Điểm GFCI chuẩn hóa của một trung tâm trong một kỳ gồm: mức riêng của trung tâm ấy ($\alpha_i$), mức chung của kỳ ấy ($\tau_t$), phần liên hệ với mức phát triển thị trường cổ phiếu của quốc gia ($\beta$), phần liên hệ với thu nhập và thể chế ($\lambda$), và phần còn lại. \HH{1} hỏi xem $\beta$ có dương không.

\begin{vidao}
\emph{Vì sao dùng điểm chuẩn hóa thay cho thứ hạng.} Thứ hạng chỉ cho biết thứ tự: khoảng cách giữa hạng 1 và hạng 2 có thể rất lớn, còn giữa hạng 60 và 61 có thể gần bằng 0, nhưng thứ hạng coi chúng như nhau. Thứ hạng cũng thay đổi khi có trung tâm mới vào bảng, như Đà Nẵng ở kỳ 40, dù các trung tâm cũ không đổi gì. Điểm z theo từng kỳ giữ được khoảng cách và đưa các kỳ về cùng thang đo, nhưng vẫn phụ thuộc vào thành phần của bảng: khi có trung tâm vào hoặc ra, trung bình và độ lệch chuẩn của kỳ thay đổi theo. Vì vậy đề tài kiểm tra độ vững trên mẫu các trung tâm có ít nhất năm kỳ và trên mẫu cân bằng (mục~\ref{sec:tincay}).
\end{vidao}

\begin{vidu}
Giả sử ước lượng được $\beta = 0{,}3$, và $\SMD$ được biểu diễn theo đơn vị độ lệch chuẩn (cách biểu diễn cụ thể sẽ được chốt trong tệp đăng ký). Khi đó, ở đặc tả có hiệu ứng cố định theo trung tâm, khi $\SMD$ của quốc gia chứa một trung tâm cao hơn mức thường lệ của chính quốc gia ấy một độ lệch chuẩn, điểm z GFCI của trung tâm ấy cao hơn mức thường lệ của nó khoảng 0,3, khi thu nhập và thể chế không đổi; chỉ ở đặc tả không có hiệu ứng cố định theo trung tâm, hệ số mới được đọc như so sánh giữa các trung tâm cùng kỳ. Đây là liên hệ có điều kiện, không phải tác động: không nói được rằng phát triển thị trường cổ phiếu \emph{làm} điểm GFCI tăng.
\end{vidu}

Sai số chuẩn được phân cụm theo quốc gia. Vì mô hình có hiệu ứng cố định theo trung tâm, $\beta$ được xác định từ biến thiên theo thời gian trong từng quốc gia, và biến thiên này nhỏ vì $\SMD$ thay đổi chậm; đề tài báo cáo thêm ước lượng không có hiệu ứng cố định theo trung tâm. Điểm GFCI đã chứa vốn hóa, giá trị và khối lượng giao dịch cổ phiếu từ Liên đoàn các Sở Giao dịch Thế giới (3 trong 144 nhân tố của kỳ GFCI 40; Wardle \& Mainelli, 2026), và trung tâm mạnh có thể kéo theo giao dịch cổ phiếu, nên $\beta$ được gọi là liên hệ có điều kiện. \HH{1} được ủng hộ khi $\beta$ dương với giá trị p nhỏ hơn ngưỡng Holm tương ứng.

\begin{delam}
\emph{Dễ nghĩ là} hiệu ứng cố định theo trung tâm luôn làm ước lượng tốt hơn. Nó loại được mọi đặc điểm không đổi của trung tâm (vị trí, ngôn ngữ, lịch sử), nhưng cái giá là $\beta$ chỉ còn dựa vào phần biến thiên nhỏ của $\SMD$ nêu ở đoạn trên. Vì vậy hai đặc tả được đọc cùng nhau.
\end{delam}

\subsubsection*{Ví dụ tính tay: từ dữ liệu ngày đến tác động}

Ví dụ dưới đây đi trọn đường của \HH{3} trên một bộ số liệu giả lập rất nhỏ: một cổ phiếu chuyển sàn (gọi là X), một cổ phiếu đối chứng HOSE (gọi là Y), mỗi mã năm ngày của một tháng trước và một tháng sau sự kiện. Mọi con số là \textbf{tự đặt để minh họa}; nghiên cứu thật dùng khoảng 20 ngày mỗi tháng, hàng trăm mã và nhiều tháng. Ngưỡng lợi suất chung trong ví dụ là 6\%.

\paragraph{Bước 1. Lợi suất tuyệt đối chia giá trị giao dịch, từng ngày.}

Với mỗi ngày, lấy lợi suất tuyệt đối (tính bằng phần trăm) chia giá trị giao dịch khớp lệnh (tính bằng tỷ đồng). Ngày có lợi suất tuyệt đối từ 6\% trở lên bị loại; ngày liền sau ngày đó cũng bị loại.

\begin{center}
{\bangnho
\begin{tabular}{@{}llrrrl@{}}
\toprule
Mã & Tháng & Ngày & Lợi suất (\%) & Giá trị (tỷ đồng) & $|r|/V$\\
\midrule
X & trước & 1 & $+1{,}0$ & 2,0 & 0,50\\
X & trước & 2 & $-8{,}0$ & 4,0 & loại (từ 6\% trở lên)\\
X & trước & 3 & $+2{,}0$ & 2,0 & loại (liền sau ngày 2)\\
X & trước & 4 & $-1{,}5$ & 1,5 & 1,00\\
X & trước & 5 & $+0{,}5$ & 1,0 & 0,50\\
\addlinespace
X & sau & 1 & $+1{,}0$ & 4,0 & 0,25\\
X & sau & 2 & $+0{,}8$ & 2,0 & 0,40\\
X & sau & 3 & $-0{,}6$ & 3,0 & 0,20\\
X & sau & 4 & $+6{,}5$ & 5,0 & loại (từ 6\% trở lên)\\
X & sau & 5 & $-1{,}2$ & 4,0 & loại (liền sau ngày 4)\\
\addlinespace
Y & trước & 1 & $+1{,}0$ & 2,5 & 0,40\\
Y & trước & 2 & $+0{,}5$ & 1,0 & 0,50\\
Y & trước & 3 & $-1{,}2$ & 2,0 & 0,60\\
Y & trước & 4 & $+0{,}9$ & 3,0 & 0,30\\
Y & trước & 5 & $-0{,}4$ & 2,0 & 0,20\\
\addlinespace
Y & sau & 1 & $+0{,}8$ & 2,0 & 0,40\\
Y & sau & 2 & $-0{,}6$ & 2,0 & 0,30\\
Y & sau & 3 & $+1{,}0$ & 4,0 & 0,25\\
Y & sau & 4 & $-0{,}3$ & 1,0 & 0,30\\
Y & sau & 5 & $+0{,}5$ & 2,0 & 0,25\\
\bottomrule
\end{tabular}}
\end{center}

\paragraph{Bước 2. Độ phi thanh khoản của tháng và logarit.}

Độ phi thanh khoản Amihud của tháng là bình quân các ngày hợp lệ; $A_{it}$ là logarit của nó.

\begin{center}
{\bangnho
\begin{tabular}{@{}llccc@{}}
\toprule
Mã & Tháng & Ngày hợp lệ & Bình quân $|r|/V$ & $A_{it}$\\
\midrule
X & trước & 1, 4, 5 & $(0{,}50 + 1{,}00 + 0{,}50)/3 = 0{,}667$ & $-0{,}405$\\
X & sau & 1, 2, 3 & $(0{,}25 + 0{,}40 + 0{,}20)/3 = 0{,}283$ & $-1{,}261$\\
Y & trước & 1--5 & $2{,}00/5 = 0{,}400$ & $-0{,}916$\\
Y & sau & 1--5 & $1{,}50/5 = 0{,}300$ & $-1{,}204$\\
\bottomrule
\end{tabular}}
\end{center}

Trong nghiên cứu thật, ở bước này còn hai việc: winsorize tại phân vị 1\% và 99\% của mẫu gộp các sàn trong tháng (mục~\ref{sec:dulieu}), và kiểm tra tháng có đủ số ngày hợp lệ tối thiểu ghi trong tệp đăng ký.

\paragraph{Bước 3. Chuỗi chênh lệch $d_t$.}

Với mỗi tháng, $d_t$ là chênh lệch logarit Amihud giữa mã chuyển sàn và mã đối chứng ghép cặp (trong nghiên cứu thật là chênh lệch trung bình trên mọi cặp):
\[
d_{\vn{trước}} = -0{,}405 - (-0{,}916) = 0{,}511 \qquad
d_{\vn{sau}} = -1{,}261 - (-1{,}204) = -0{,}057
\]
Trước sự kiện, X kém thanh khoản hơn Y ($d$ dương). Sau sự kiện, X gần bằng Y.

\paragraph{Bước 4. Thống kê kiểm định và tác động.}

Thống kê kiểm định là trung bình $d_t$ sau sự kiện trừ trung bình $d_t$ trong cửa sổ trước. Ở đây mỗi giai đoạn chỉ có một tháng:
\[
-0{,}057 - 0{,}511 = -0{,}568
\]
Quy ra phần trăm: $e^{-0{,}568} - 1 \approx -0{,}43$, tức độ phi thanh khoản của X giảm khoảng 43\% so với Y. Đây cũng là $\delta$ của phương trình \eqref{eq:did} khi chỉ có hai mã và hai tháng.

\paragraph{Bước 5. Giá trị p.}

Với một tháng mỗi phía thì không tính được giá trị p. Trong nghiên cứu thật, chuỗi $d_t$ có hàng chục tháng trước và 24 tháng sau; mô hình tự hồi quy được ước lượng trên giai đoạn tham chiếu và bootstrap tự hồi quy cho biết một giá trị như $-0{,}568$ có hiếm gặp khi không có sự kiện hay không (mục~\ref{sec:kiemdinh}).

\paragraph{Điều ví dụ này cho thấy.}

\begin{itemize}
\item Ngưỡng chung làm X mất hai ngày ở mỗi tháng, ở cả trước lẫn sau, theo cùng một quy tắc. Nếu không có ngưỡng, ngày $-8\%$ ở HNX vẫn được tính, còn một ngày tương tự ở HOSE sẽ bị chặn ở 7\%; hai tháng sẽ được tính trên hai cơ sở khác nhau.
\item Y cũng tốt lên ($-1{,}204$ so với $-0{,}916$): nếu chỉ nhìn X trước và sau, ta sẽ gán cả phần cải thiện chung của thị trường cho việc chuyển sàn.
\item Kết quả từ hai mã và mười ngày không nói lên điều gì về thực tế; ví dụ này chỉ để thấy từng bước tính.
\end{itemize}


%-------------------------------------------------------
\subsection{Thiết kế khác biệt kép quanh ngày chuyển sàn}\label{sec:did}
%-------------------------------------------------------

\begin{tomtat}
Khác biệt kép chỉ đúng khi ba giả định đúng: nếu không có sự kiện thì hai nhóm đi song song; nhóm đối chứng không bị sự kiện tác động; và việc cổ phiếu rời mẫu không liên quan đến sự kiện. Mục này nói cách kiểm tra từng giả định, nhất là các thay đổi trùng thời gian ở HOSE (FTSE Russell, phân bảng, VN-Index) và khả năng thanh khoản lan giữa hai nhóm.
\end{tomtat}

Đề tài kiểm tra ba giả định của thiết kế khác biệt kép. Thứ nhất là xu hướng song song, được kiểm tra bằng các hệ số trước sự kiện của phương trình \eqref{eq:es}.

Thứ hai, nhóm đối chứng không chịu tác động của sự kiện. Nhóm đối chứng HOSE chịu ba thay đổi trùng thời gian: FTSE Russell đưa cổ phiếu Việt Nam vào chỉ số theo bốn đợt từ tháng 9/2026 đến tháng 9/2027 (FTSE Russell, 2026b); HOSE dự kiến công bố phân loại cổ phiếu vào Bảng Chính và Bảng Tiêu chuẩn lần đầu ngày 28/12/2026, theo dự thảo quy chế đang lấy ý kiến (Tuổi Trẻ Online, 2026a); và VN-Index, vốn tự động cập nhật danh mục sau ngày giao dịch đầu tiên của cổ phiếu niêm yết mới (Sở Giao dịch Chứng khoán Thành phố Hồ Chí Minh, 2026), nhiều khả năng nhận các mã chuyển sàn từ phiên thứ hai trên HOSE. Ngày vào chỉ số thực tế của từng mã và quy chế phân bảng chính thức, khi được ban hành, được ghi lại từ công bố của HOSE và đưa vào ước lượng.

Vì vậy, nhóm đối chứng HOSE loại các mã thuộc danh sách đủ điều kiện của FTSE Russell, $\delta$ được ước lượng thêm với biến giả nhóm bảng phân loại, và thanh khoản của từng nhóm đối chứng được so với xu hướng trước của chính nhóm đó.

\begin{center}
\captionof{table}{Thay đổi trùng thời gian ở HOSE và cách xử lý}\label{tab:trungthoigian}
{\bangnho
\begin{tabularx}{\textwidth}{@{}>{\raggedright\arraybackslash}p{0.27\textwidth} L L@{}}
\toprule
Thay đổi trùng thời gian & Có thể ảnh hưởng & Cách xử lý\\
\midrule
FTSE Russell đưa cổ phiếu Việt Nam vào chỉ số (bốn đợt, 9/2026--9/2027) & Thanh khoản các mã HOSE đủ điều kiện tăng & Loại các mã đủ điều kiện khỏi nhóm đối chứng; nếu mã chuyển sàn được đưa vào, ước lượng thêm khi loại các mã này\\
Phân loại Bảng Chính và Bảng Tiêu chuẩn (dự kiến 28/12/2026) & Hai nhóm bảng có thể được nhà đầu tư đối xử khác nhau & Ước lượng thêm với biến giả nhóm bảng\\
VN-Index nhận mã chuyển sàn & Áp lực cầu lên mã chuyển sàn, nhất là mã có vốn hóa lớn so với giá trị giao dịch & Đo bằng mức chịu cú sốc chỉ số (mục~\ref{sec:kiemdinh})\\
Nới biên độ, giao dịch xuyên trưa (đang nghiên cứu) & Mọi mã HOSE & Ghi ngày áp dụng; báo cáo thêm ước lượng với cửa sổ sau cắt tại ngày đó\\
\bottomrule
\end{tabularx}}
\end{center}

Ủy ban Chứng khoán Nhà nước cũng đang nghiên cứu nới biên độ và giao dịch xuyên trưa (Thời báo Tài chính Việt Nam, 2026). Nếu các thay đổi này được áp dụng trong cửa sổ sau, ngày áp dụng được ghi là một cú sốc trùng thời gian và ước lượng được báo cáo thêm với cửa sổ sau cắt tại ngày đó.

Nếu nhóm đối chứng HOSE mất thanh khoản do dòng lệnh chuyển sang cổ phiếu chuyển sàn, $\delta$ phóng đại tác động, và ước lượng với nhóm đối chứng UPCoM, nếu nhóm này còn dữ liệu sau sự kiện, là một cận khác. Ngược lại, nếu thanh khoản lan tỏa dương sang nhóm đối chứng HOSE, như đã ghi nhận ở Tel Aviv (Amihud và cộng sự, 1997) và quanh đợt hợp nhất Tokyo--Osaka (Hisada, 2018), $\delta$ đánh giá thấp tác động.

\begin{vidu}[: hai hướng lan tỏa]
Giả sử tác động thật của chuyển sàn là $-0{,}20$. Nếu nhà đầu tư rút tiền khỏi các mã HOSE đối chứng để mua mã chuyển sàn, nhóm đối chứng kém thanh khoản đi, chẳng hạn $+0{,}05$. Khác biệt kép đo được là $-0{,}20 - 0{,}05 = -0{,}25$, phóng đại tác động. Nếu ngược lại, sàn HOSE đông giao dịch hơn làm mọi mã cùng dễ mua bán, nhóm đối chứng tốt lên (logarit Amihud giảm 0,05, tức thay đổi $-0{,}05$), và khác biệt kép đo được là $-0{,}15$, đánh giá thấp tác động. Vì không biết trước hướng nào mạnh hơn, đề tài dùng nhóm đối chứng UPCoM, ít liên quan đến dòng lệnh trên HOSE, làm một cận khác.
\end{vidu}

Thứ ba, việc rời mẫu không liên quan đến sự kiện; đề tài báo cáo tỷ lệ mã hủy niêm yết ở từng nhóm và ước lượng lại trên bảng cân bằng. Vì việc chuyển sàn đổi cùng lúc biên độ, bước giá, hệ thống giao dịch và tư cách thành phần chỉ số, $\delta$ đo tác động gộp của sàn, không đo riêng từng quy tắc.

%-------------------------------------------------------
\subsection{Quy trình nghiên cứu}\label{sec:quytrinh}
%-------------------------------------------------------

Hình~\ref{fig:quytrinh} tóm tắt quy trình gồm bước đăng ký trước (bước 0) và năm bước nghiên cứu: thu thập dữ liệu quốc tế và dữ liệu ba sàn; xây dựng Chỉ số chất lượng thị trường vốn và kiểm định \HH{2}; kiểm định \HH{1} trên dữ liệu bảng quốc tế; ước lượng tác động của sự kiện chuyển sàn lên thanh khoản (\HH{3}) và khả năng đầu tư (\HH{4}); tổng hợp định hướng chuyên môn hóa cho thị trường vốn Hà Nội. Thí sinh đăng ký trước kế hoạch phân tích trước ngày 28/12/2026.

\begin{figure}[H]
\myfloatalign
\begin{tikzpicture}[node distance=4mm,
  buoc/.style={draw, rectangle, align=left, text width=0.78\textwidth, minimum height=9mm, font=\footnotesize, inner sep=4pt},
  >=stealth]
\node[buoc] (b0) {\textbf{Bước 0.} Đăng ký trước kế hoạch phân tích trên OSF Registries (trước 28/12/2026)};
\node[buoc, below=of b0] (b1) {\textbf{Bước 1.} Thu thập dữ liệu: điểm GFCI kỳ 1--40, chỉ báo quốc gia; dữ liệu ngày của HNX, UPCoM, HOSE từ 6/2009};
\node[buoc, below=of b1] (b2) {\textbf{Bước 2.} Xây dựng Chỉ số chất lượng thị trường vốn; kiểm định \HH{2}};
\node[buoc, below=of b2] (b3) {\textbf{Bước 3.} Kiểm định \HH{1} trên dữ liệu bảng quốc tế};
\node[buoc, below=of b3] (b4) {\textbf{Bước 4.} Khác biệt kép quanh 28/12/2026: thanh khoản (\HH{3}), khả năng đầu tư (\HH{4})};
\node[buoc, below=of b4] (b5) {\textbf{Bước 5.} Tổng hợp định hướng chuyên môn hóa cho thị trường vốn Hà Nội};
\foreach \a/\b in {b0/b1,b1/b2,b2/b3,b3/b4,b4/b5}{\draw[->] (\a) -- (\b);}
\end{tikzpicture}
\caption{Quy trình nghiên cứu}\label{fig:quytrinh}
\nguon{Nguồn: Thí sinh đề xuất. Ghi chú: bước 0 là điều kiện của mọi bước sau; bước 2 và bước 3 có thể làm song song.}
\end{figure}

\paragraph{Từ bằng chứng đến định hướng cho Hà Nội.} Kết quả \HH{1}--\HH{4} được tổng hợp cùng phép so sánh vai trò của Hà Nội với Bắc Kinh và Thượng Hải (Lai, 2012; Wang và cộng sự, 2007) và với Thành phố Hồ Chí Minh, dựa trên văn bản chính sách và dữ liệu công khai, thành định hướng chuyên môn hóa cho các thị trường mà HNX tiếp tục vận hành. Phần này là tổng hợp định tính, không kiểm định giả thuyết.

Cách tổng hợp có thể hình dung qua bốn tình huống. Nếu chuyển sàn làm thanh khoản và khả năng đầu tư tăng rõ, thì việc chuyển sàn có lợi cho nhà đầu tư vào các mã chuyển sàn, và thế mạnh của Hà Nội nên được tìm ở các thị trường khác ngoài cổ phiếu niêm yết. Nếu khoảng tin cậy hẹp quanh 0, khác biệt giữa hai sàn ở \HH{2} nhiều khả năng đến từ đặc điểm doanh nghiệp hơn là từ sàn; nếu kết quả chưa đủ cơ sở kết luận, đề tài không rút ra hàm ý nào về vai trò của sàn. Nếu doanh nghiệp Hà Nội được lợi ít hơn, thì có lý do để thành phố quan tâm tới hạ tầng thông tin cho doanh nghiệp địa phương. Nếu \HH{1} được ủng hộ, thì ở cấp quốc gia, điểm GFCI đi cùng mức phát triển thị trường cổ phiếu (liên hệ có điều kiện). Vì $\SMD$ đo ở cấp quốc gia và như nhau với mọi thành phố của Việt Nam, kết quả này không cho biết việc giao dịch cổ phiếu tập trung ở Thành phố Hồ Chí Minh có bất lợi cho Hà Nội hay không; nó chỉ là bối cảnh cho phần so sánh định tính. Bốn tình huống này chỉ là khung đọc kết quả; đề tài không chọn trước tình huống nào.

